\documentclass{beamer}

\usepackage[T1]{fontenc}
\usepackage[utf8]{inputenc}
\usepackage[frenchb]{babel}
\usepackage{eurosym}

\usepackage{multirow}

\usepackage{graphicx}


\newcommand\host{{\it host}\,}
\newcommand\device{{\it device}\,}


% --------- BEGIN STYLE BEAMER ---------
\usetheme{Madrid}
\definecolor{ECNBleu}{HTML}{102648}
\definecolor{ECNJaune}{HTML}{FAB600}

\setbeamercolor{frametitle}{bg=ECNBleu,fg=white} %le titre de la frame, en haut
\setbeamercolor{title in head/foot}{bg=ECNJaune,fg=ECNBleu} %barre bas, milieu
\setbeamercolor{author in head/foot}{bg=ECNBleu,fg=ECNJaune} %barre bas, gauche
\setbeamercolor{date in head/foot}{bg=ECNBleu,fg=ECNJaune} %barre bas, droit
\setbeamercolor{section in head/foot}{bg=ECNJaune,fg=ECNBleu} %haut gauche
\setbeamercolor{subsection in head/foot}{bg=ECNBleu,fg=ECNBleu} %haut droite
\setbeamercolor{title}{bg=ECNBleu,fg=white}
\setbeamercolor{block title}{fg=white,bg=ECNBleu}

%\setbeamercolor{normal text}{bg=IUTBleuFonce} %texte
\setbeamercolor{item}{bg=ECNBleu}
\setbeamercolor{subitem}{bg=ECNJaune}
\setbeamercolor{description item}{fg=couleurEmph}
\setbeamercolor{section in toc}{fg=ECNBleu}
%\setbeamercolor{subsubitem}{bg=IUTBleuFonce}

%structure (eq emph dans beamer)
\setbeamercolor{emph}{fg=couleurEmph}
% ----------------------------

\usepackage{listings}

\lstdefinestyle{customc}{
  belowcaptionskip=1\baselineskip,
  breaklines=true,
  frame=L,
  xleftmargin=\parindent,
  language=C,
  showstringspaces=false,
  basicstyle=\tiny\ttfamily,
  keywordstyle=\bfseries\color{green!40!black},
  commentstyle=\itshape\color{purple!40!black},
  identifierstyle=\color{blue},
  stringstyle=\color{orange},
}

\lstset{escapechar=@,style=customc}

\newenvironment<>{questionblock}[1]{%
\setbeamercolor{block title}{bg=ECNJaune,fg=ECNBleu}%
\begin{block}#2{#1}}{\end{block}}

\title[GPGPU]{Programmation sur Processeur Graphique}

\subtitle[\ldots]{
---\\Partie 5 : détour sur la virgule flottante}
\author[P.-E. Hladik]{Pierre-Emmanuel Hladik\\\tiny pierre-emmanuel.hladik@ec-nantes.fr}
\institute[Centrale Nantes]{}
\date{\today}

\begin{document}

 \frame{\titlepage}

%%%%%%%%%%%%%%%%%%%%%%%%%%%
%\begin{frame}{Plan}
%\tableofcontents[sections={1-4}, currentsection, hideothersubsections]
%\end{frame}
%%%%%%%%%%%%%%%%%%%%%%%%%%%

%%%%%%%%%%%%%%%%%%%%%%%%%%
\begin{frame}[plain]
\frametitle{Qu'est que le format à virgule flottante ?}

	\begin{itemize}
		\item Une norme industrielle (IEEE 754 2008) pour représenter les nombres à virgule flottante afin de garantir que les matériels de différents fournisseurs génèrent des résultats cohérents entre eux\\
		\item Un nombre binaire à virgule flottante se compose de trois parties :
		\begin{itemize}
			\item le signe (S), l'exposant (E) et la mantisse (M) [entier relatif, généralement borné]
			\item Chaque tuple (S, E, M) identifie de manière unique un nombre à virgule flottante\\
		\end{itemize}
		\item Pour chaque motif binaire, sa valeur en virgule flottante IEEE est dérivée comme suit :
		\begin{itemize}
			\item valeur = $(-1)^S \cdot M \cdot 2^E$ avec $1.0 \leq M < 10.0_B$
		\end{itemize}
		\item L'interprétation de S est simple : $S=0$ pour un nombre positif et $S=1$ à un nombre négatif
	\end{itemize}

\end{frame}

%%%%%%%%%%%%%%%%%%%%%%%%%%
\begin{frame}[plain]
\frametitle{Représentation normalisée}

\begin{block}{}
	En spécifiant que $1.0B \leq M<10.0_B$, on rend unique la valeur de la mantisse pour chaque nombre à virgule flottante
\end{block}
		\begin{itemize}
			\item Par exemple, pour $0.5_D$ la seule valeur de mantisse autorisée est $M =1.0$ :  $0.5_D  =1.0_B \cdot 2^{-1}$
			\item ça ne peut pas être $10.0_B \cdot 2^{ -2}$, ni $0.1_B \cdot 2^{0}$
		\end{itemize}

\begin{block}{}
	Comme toutes les valeurs de mantisse sont de la forme $1.XX...$ on peut omettre la partie "$1.$" dans la représentation (on parle de bit implicite)
\end{block}

		\begin{itemize}
			\item Par exemple, la mantisse de $0.5_D$ dans une mantisse de 2 bits est $00$, qui est obtenue en omettant "1." de $1.00_B$.
		\end{itemize}

\end{frame}

%%%%%%%%%%%%%%%%%%%%%%%%%%
\begin{frame}[plain]
\frametitle{Représentation de l'exposant (IEEE) (1/3)}

\begin{block}{Complément à deux}
Méthode pour représenter des entiers relatifs en binaire pour effectuer simplement des opérations arithmétiques.

Les nombres négatifs sont obtenus en calculant l'opposé du nombre positif par deux opérations successives:
\begin{itemize}
		\item On inverse les bits de l'écriture binaire (opération binaire NON), on fait ce qu'on appelle le complément à un ;
		\item On ajoute 1 au résultat (les dépassements sont ignorés).
	\end{itemize}
\end{block}

Exemple : codage binaire de $-4$ sur 8 bits :
\begin{itemize}
		\item codage du nombre positif 4 : $0000 \ 0100_B$
	\item inversion des bits : $1111 \ 1011_B$
	\item ajout de 1 : $1111 \ 1100_B$
\end{itemize}

\end{frame}
%%%%%%%%%%%%%%%%%%%%%%%%%%

%%%%%%%%%%%%%%%%%%%%%%%%%%
\begin{frame}[plain]
\frametitle{Représentation de l'exposant (IEEE) (2/3)}


\begin{block}{Propriétés}
	\begin{itemize}
		\item -0 est 0 : $0000 \ 0000_B$ devient $1111 \ 1111_B + 1$ soit $0000 \ 0000_B$
		\item L'addition usuelle des nombres binaires fonctionne : $3 + (-4) = 0011_B + 1100_B = 1111_B = -1$
\end{itemize}
\end{block}


\end{frame}
%%%%%%%%%%%%%%%%%%%%%%%%%%

%%%%%%%%%%%%%%%%%%%%%%%%%%
\begin{frame}[plain]
\frametitle{Représentation de l'exposant (IEEE) (3/3)}

	\begin{itemize}
		\item Pour un exposant $E$ de $n$ bits on ajoute $2^{n-1}-1$à la représentation de son complément à 2  (exemple pour une représentation de l'exposant sur 3 bits)
	\item En général, avec une mantisse normalisée et un exposant codé en excès, la valeur d'un nombre avec un exposant de n bits est :
$(-1)S \cdot 1.M \cdot 2^{(E- (2^{(n - 1)} - 1))}$
	\end{itemize}

\begin{table}[htdp]
\begin{center}
\begin{tabular}{|c|c|c|}
\hline
$E$ & Valeur décimale & $Excess-3$ \\
\hline
101 & –3& 000\\
110 &–2 &001\\
111 &–1 &010\\
000 &0 &011\\
001  &1 &100\\
010  &2 &101\\
011  &3& 110\\
100 &Reserved pattern &111\\
\hline
\end{tabular}
\end{center}
\end{table}%

\end{frame}
%%%%%%%%%%%%%%%%%%%%%%%%%%

\begin{frame}[plain]
\frametitle{Exemple sur un format 5-bit flottant}

Supposons un format avec 1bit pour S, 2 bits pour E et 2 bits pour M, la valeur $0.5_D = 1.00_B \cdot 2^{-1} =  0 \ 00 \ 00$ avec $S = 0$, $E = 00$ et $M = (1.)00$

\begin{table}[htdp]
\begin{center}
\begin{tabular}{|c|c|c|}
\hline
$E$ & Valeur décimale & $Excess-1$ \\
\hline
11 &–1 &00\\
00 &0 &01\\
01  &1 &11\\
10 &Reserved pattern &11\\
\hline
\end{tabular}
\end{center}
\end{table}%

\end{frame}

%%%%%%%%%%%%%%%%%%%%%%%%%%
\begin{frame}[plain]
\frametitle{Nombres représentables}

	\begin{itemize}
		\item Les nombres représentables d'un format donné sont l'ensemble de tous les nombres qui peuvent être représentés exactement dans ce format.
		\item Voir le tableau des nombres représentables d'un format d'entier de 3 bits non signé.
	\end{itemize}
	
\begin{table}[htdp]
\begin{center}
\begin{tabular}{|c|c|}
\hline
000 & 0\\
001 & 1\\
010& 2\\
011& 3\\
100& 4\\
101& 5\\
110& 6\\
111& 7\\
\hline
\end{tabular}
\end{center}
\end{table}%

\end{frame}

%%%%%%%%%%%%%%%%%%%%%%%%%%
\begin{frame}[plain]
\frametitle{Nombres représentables en 5 bits}


\begin{table}[htdp]
\begin{center}
\begin{tabular}{|c|c|c|c|}
\hline
E & M & S = 0 & S = 1\\
\hline
\multirow{4}{*}{00}	& 00 & $2^{-1}$ & $- (2^{-1})$\\
				& 01 & $2^{-1} + 1\cdot 2^{-3}$ & $- (2^{-1} + 1\cdot 2^{-3})$\\
				& 10 & $2^{-1} + 2\cdot 2^{-3}$ & $- (2^{-1} + 2\cdot 2^{-3})$\\
				& 11 & $2^{-1} + 3\cdot 2^{-3}$ & $- (2^{-1} + 3\cdot 2^{-3})$\\
\hline
\multirow{4}{*}{01}	& 00 & $2^{0}$ & $- (2^{0})$\\
				& 01 & $2^{0} + 1\cdot 2^{-3}$ & $- (2^{0} + 1\cdot 2^{-3})$\\
				& 10 & $2^{0} + 2\cdot 2^{-3}$ & $- (2^{0} + 2\cdot 2^{-3})$\\
				& 11 & $2^{0} + 3\cdot 2^{-3}$ & $- (2^{0} + 3\cdot 2^{-3})$\\
\hline
\multirow{4}{*}{10}	& 00 & $2^{1}$ & $- (2^{1})$\\
				& 01 & $2^{1} + 1\cdot 2^{-3}$ & $- (2^{1} + 1\cdot 2^{-3})$\\
				& 10 & $2^{1} + 2\cdot 2^{-3}$ & $- (2^{1} + 2\cdot 2^{-3})$\\
				& 11 & $2^{1} + 3\cdot 2^{-3}$ & $- (2^{1} + 3\cdot 2^{-3})$\\
\hline
11 & \multicolumn{3}{c|}{Réservé}\\
\hline
\end{tabular}
\end{center}
\end{table}%

\end{frame}
%%%%%%%%%%%%%%%%%%%%%%%%%%
\begin{frame}[plain]
\frametitle{Observations sur la représentation en 5 bits}

\begin{itemize}
		\item les bits de l'exposant définissent les intervalles majeurs des nombres représentables. Les intervalles majeurs sont entre des puissances de 2
		\item Les bits de la mantisse définissent le nombre de nombres représentatifs dans chaque intervalle majeur. Par conséquent, le nombre de bits de la mantisse détermine la précision de la représentation
		\item Les nombres représentables se rapprochent les uns des autres dans le voisinage de 0. C'est une tendance souhaitable, car plus la valeur absolue de ces nombres est faible, plus il est important de les représenter avec précision.
		\item la tendance à l'augmentation de la densité des nombres représentables, et donc à l'augmentation de la précision de la représentation des nombres dans les intervalles au fur et à mesure que l'on se rapproche de zéro, ne s'applique pas aux environs immédiats de 0. Cela est dû au fait que la plage de mantisse normalisée exclut 0. Il s'agit d'une autre déficience grave.
	\end{itemize}

\end{frame}

%%%%%%%%%%%%%%%%%%%%%%%%%%
\begin{frame}[plain]
\frametitle{Mise à zéro}

Traiter toutes les configurations binaires avec $E=0$ comme étant $0.0$

\begin{itemize}
		\item Cela enlève plusieurs nombres représentables proches de zéro et les regroupe tous dans 0.0
		\item Pour une représentation avec un grand M, un grand nombre de nombres représentables sont supprimés.
	\end{itemize}
	
	\begin{table}[htdp]
	
\tiny 
\begin{center}
\begin{tabular}{|c|c|c|c||c|c|}
\cline{3-6}
\multicolumn{2}{c|}{ } & \multicolumn{2}{c|}{pas de zéro} & \multicolumn{2}{c|}{mise à zéro}\\
\hline
E & M & S = 0 & S = 1& S = 0 & S = 1\\
\hline
\multirow{4}{*}{00}	& 00 & $2^{-1}$ & $- (2^{-1})$ & 0 & 0\\
				& 01 & $2^{-1} + 1\cdot 2^{-3}$ & $- (2^{-1} + 1\cdot 2^{-3})$& 0 & 0\\
				& 10 & $2^{-1} + 2\cdot 2^{-3}$ & $- (2^{-1} + 2\cdot 2^{-3})$& 0 & 0\\
				& 11 & $2^{-1} + 3\cdot 2^{-3}$ & $- (2^{-1} + 3\cdot 2^{-3})$& 0 & 0\\
\hline
\multirow{4}{*}{01}	& 00 & $2^{0}$ & $- (2^{0})$& $2^{0}$ & $- (2^{0})$\\
				& 01 & $2^{0} + 1\cdot 2^{-3}$ & $- (2^{0} + 1\cdot 2^{-3})$ & $2^{0} + 1\cdot 2^{-3}$ & $- (2^{0} + 1\cdot 2^{-3})$\\
				& 10 & $2^{0} + 2\cdot 2^{-3}$ & $- (2^{0} + 2\cdot 2^{-3})$ & $2^{0} + 2\cdot 2^{-3}$ & $- (2^{0} + 2\cdot 2^{-3})$\\
				& 11 & $2^{0} + 3\cdot 2^{-3}$ & $- (2^{0} + 3\cdot 2^{-3})$ & $2^{0} + 3\cdot 2^{-3}$ & $- (2^{0} + 3\cdot 2^{-3})$\\
\hline
\multirow{4}{*}{10}	& 00 & $2^{1}$ & $- (2^{1})$ & $2^{1}$ & $- (2^{1})$\\
				& 01 & $2^{1} + 1\cdot 2^{-3}$ & $- (2^{1} + 1\cdot 2^{-3})$ & $2^{1} + 1\cdot 2^{-3}$ & $- (2^{1} + 1\cdot 2^{-3})$\\
				& 10 & $2^{1} + 2\cdot 2^{-3}$ & $- (2^{1} + 2\cdot 2^{-3})$ & $2^{1} + 2\cdot 2^{-3}$ & $- (2^{1} + 2\cdot 2^{-3})$\\
				& 11 & $2^{1} + 3\cdot 2^{-3}$ & $- (2^{1} + 3\cdot 2^{-3})$ & $2^{1} + 3\cdot 2^{-3}$ & $- (2^{1} + 3\cdot 2^{-3})$\\
\hline
11 & \multicolumn{5}{c|}{Réservé}\\
\hline
\end{tabular}
\end{center}
\end{table}%

\end{frame}

%%%%%%%%%%%%%%%%%%%%%%%%%%

%%%%%%%%%%%%%%%%%%%%%%%%%%
\begin{frame}[plain]
\frametitle{Nombres dénormalisés (format IEEE)}

\begin{itemize}
		\item La méthode assouplit l'exigence de normalisation pour les nombres très proches de 0
		\item Lorsque $E=0$, la mantisse n'est plus supposée être de la forme $1.XX$ mais $0.XX$, la valeur est $0.M * 2^{- 2 ^{(n-1)} + 2}$
	\end{itemize}
	
	\begin{table}[htdp]
	
\tiny 
\begin{center}
\begin{tabular}{|c|c|c|c||c|c|}
\cline{3-6}
\multicolumn{2}{c|}{ } & \multicolumn{2}{c|}{pas de zéro} & \multicolumn{2}{c|}{Dénormalisé}\\
\hline
E & M & S = 0 & S = 1& S = 0 & S = 1\\
\hline
\multirow{4}{*}{00}	& 00 & $2^{-1}$ & $- (2^{-1})$ & 0 & 0\\
				& 01 & $2^{-1} + 1\cdot 2^{-3}$ & $- (2^{-1} + 1\cdot 2^{-3})$& $1\cdot 2^{-2}$ & $- 1\cdot 2^{-2}$\\
				& 10 & $2^{-1} + 2\cdot 2^{-3}$ & $- (2^{-1} + 2\cdot 2^{-3})$& $2\cdot 2^{-2}$ & $- 2\cdot 2^{-2}$\\
				& 11 & $2^{-1} + 3\cdot 2^{-3}$ & $- (2^{-1} + 3\cdot 2^{-3})$& $3\cdot 2^{-2}$ & $- 3\cdot 2^{-2}$\\
\hline
\multirow{4}{*}{01}	& 00 & $2^{0}$ & $- (2^{0})$& $2^{0}$ & $- (2^{0})$\\
				& 01 & $2^{0} + 1\cdot 2^{-3}$ & $- (2^{0} + 1\cdot 2^{-3})$ & $2^{0} + 1\cdot 2^{-3}$ & $- (2^{0} + 1\cdot 2^{-3})$\\
				& 10 & $2^{0} + 2\cdot 2^{-3}$ & $- (2^{0} + 2\cdot 2^{-3})$ & $2^{0} + 2\cdot 2^{-3}$ & $- (2^{0} + 2\cdot 2^{-3})$\\
				& 11 & $2^{0} + 3\cdot 2^{-3}$ & $- (2^{0} + 3\cdot 2^{-3})$ & $2^{0} + 3\cdot 2^{-3}$ & $- (2^{0} + 3\cdot 2^{-3})$\\
\hline
\multirow{4}{*}{10}	& 00 & $2^{1}$ & $- (2^{1})$ & $2^{1}$ & $- (2^{1})$\\
				& 01 & $2^{1} + 1\cdot 2^{-3}$ & $- (2^{1} + 1\cdot 2^{-3})$ & $2^{1} + 1\cdot 2^{-3}$ & $- (2^{1} + 1\cdot 2^{-3})$\\
				& 10 & $2^{1} + 2\cdot 2^{-3}$ & $- (2^{1} + 2\cdot 2^{-3})$ & $2^{1} + 2\cdot 2^{-3}$ & $- (2^{1} + 2\cdot 2^{-3})$\\
				& 11 & $2^{1} + 3\cdot 2^{-3}$ & $- (2^{1} + 3\cdot 2^{-3})$ & $2^{1} + 3\cdot 2^{-3}$ & $- (2^{1} + 3\cdot 2^{-3})$\\
\hline
11 & \multicolumn{5}{c|}{Réservé}\\
\hline
\end{tabular}
\end{center}
\end{table}%

\end{frame}

%%%%%%%%%%%%%%%%%%%%%%%%%%

%%%%%%%%%%%%%%%%%%%%%%%%%%
\begin{frame}[plain]
\frametitle{Format IEEE et précision}

\begin{itemize}
		\item Précision simple (32 bits) :  Signe 1 bit, exposant 8 bits (excès de biais-127), mantisse 23 bits
		\item Double précision (64 bits) : Signe de 1 bit, exposant de 11 bits (excès de biais de 1023), mantisse de 52 bits.
		\item En double précision, l'erreur la plus importante pour représenter un nombre est réduite à 1/229 de la représentation en simple précision.
\end{itemize}
	
	\begin{table}[htdp]
\caption{Patterns spéciaux}
\begin{center}
\begin{tabular}{|c|c|c|}
\hline
Exposant & Mantisse & Signification\\
\hline
11...1 & $\neq 0$ & NaN (Not a Number)\\
11...1 & $= 0$ & $(-1)^S \infty$\\
00...0 & $\neq 0$ & dénormalisé\\
00...0 & $= 0$ & 0\\
\hline
\end{tabular}
\end{center}
\label{default}
\end{table}%

	
\end{frame}

%%%%%%%%%%%%%%%%%%%%%%%%%%

\begin{frame}[plain]
\frametitle{Précision et arrondi en virgule flottante}

\begin{itemize}
		\item La précision d'une opération arithmétique à virgule flottante est mesurée par l'erreur maximale introduite par l'opération
		\item La source d'erreur la plus courante dans l'arithmétique à virgule flottante est le fait que l'opération génère un résultat qui ne peut pas être représenté exactement et qui doit donc être arrondi
		\item L'arrondi se produit lorsque la mantisse de la valeur doit être représentée sur trop de bits pour être exacte
\end{itemize}


\end{frame}

%%%%%%%%%%%%%%%%%%%%%%%%%%

\begin{frame}[plain]
\frametitle{Arrondi et erreur}

Considérons une représentation 5-bits et l'addition de $1.00_B \cdot 2^{-2} $ avec $1.00_B \cdot 2^{1}$, soit (0,00,01) +  (0,10,00).

Le matériel doit décaler les bits de la mantisse pour les aligner, soit $0.001_B \cdot 2^{1} + 1.00_B \cdot 2^{1} = 1.001_B \cdot 2^{1}$... qui n'est pas représentable avec une mantisse de 2 bits.

La représentation la plus proche est $1.00_B \cdot 2^{1}$ ou $1.01_B \cdot 2^{1}$ soit une erreur de $0.001_B\cdot 2^1$ qui est la moitié de la valeur de la place la moins significative de la mantisse.


\end{frame}

%%%%%%%%%%%%%%%%%%%%%%%%%%

\begin{frame}[plain]
\frametitle{L'ordre des opérations peut être important}

\begin{itemize}
		\item Les opérations en virgule flottante ne sont pas strictement associatives.
		\item La cause première est que, parfois, un nombre petit peut disparaître lorsqu'il est ajouté ou soustrait à un nombre très grand, soit :
		
		  (Grand + Petit) + Petit $\neq$ Grand + (Petit + Petit)
\end{itemize}

Exemple : $1.00_B \cdot 2^{0} +1.00_B \cdot 2^{0} + 1.00_B \cdot 2^{-2} + 1.00_B \cdot 2^{-2}$

\begin{itemize}
\item En séquentiel : \\
$1.00_B \cdot 2^{0} +1.00_B \cdot 2^{0} + 1.00_B \cdot 2^{-2} + 1.00_B \cdot 2^{-2} $\\
$= 1.00_B \cdot 2^{1} + 1.00_B \cdot 2^{-2} + 1.00_B \cdot 2^{-2}$\\
$= 1.00_B \cdot 2^{1} + 1.00_B \cdot 2^{-2}$ (arrondi)\\
$= 1.00_B \cdot 2^{1}$ (arrondi)
\item En parallèle : \\
$(1.00_B \cdot 2^{0} +1.00_B \cdot 2^{0}) + (1.00_B \cdot 2^{-2} + 1.00_B \cdot 2^{-2}) $\\
$= 1.00_B \cdot 2^{1} + 1.00_B \cdot 2^{-1} $\\
$= 1.01_B \cdot 2^{1}$

\end{itemize}


\end{frame}
%%%%%%%%%%%%%%%%%%%%%%%%%%

\begin{frame}[plain]
\frametitle{Technique pour maximiser la précision}

\begin{block}{}
Une technique pour maximiser la précision arithmétique en virgule flottante consiste à pré-trier les données avant un calcul de réduction.
\end{block}

Dans notre exemple de réduction de la somme, si nous pré-trions les données par ordre numérique croissant, nous obtiendrons ce qui suit :
$ 1.00_B \cdot 2^{-2} + 1.00_B \cdot 2^{-2}  + 1.00_B \cdot 2^{0} +1.00_B \cdot 2^{0}$\\
$= 1.00_B \cdot 2^{-1} + 1.00_B \cdot 2^{0} +1.00_B \cdot 2^{0}$\\
$= 1.10_B \cdot 2^{0} +1.00_B \cdot 2^{0}$\\
$= 1.01_B \cdot 2^{1}$

\end{frame}

%More specifically, all CUDA devices of compute capabil- ity 1.3 and up support denormalized double-precision operands, and all devices of compute capability 2.0 and up support denormalized single-precision operands.

%%%%%%%%%%%%%%%%%%%%%%%%%%

\end{document}